% Original teaching examples; numerical features are illustrative.
\begin{tikzpicture}[x=1mm,y=1mm,font=\figurefont\footnotesize,
  every node/.style={text=NNInk},
  nn flow/.append style={draw=NNWire,rounded corners=1.5mm},
  nn operator path/.append style={draw=NNWire,rounded corners=1.5mm},
  token/.style={draw=NNInput,fill=NNInput!18,line width=.85pt,minimum width=9mm,minimum height=7mm,inner sep=1pt},
  module/.style={nn operator,inner sep=3pt}]
  \node[nn heading,anchor=west] at (0,2) {(a) 稀疏：只保留邻近位置对};
  \foreach \i/\word in {1/猫,2/在,3/院,4/里,5/追,6/狗}{
    \node at ({14+6*\i},-8) {\word};
    \node at (12,{-11-6*\i}) {\word};
    \foreach \j in {1,...,6}{
      \pgfmathtruncatemacro{\keep}{abs(\i-\j)<=1}
      \ifnum\keep=1
        \fill[NNHidden!25] ({11+6*\j},{-14-6*\i}) rectangle ++(6,6);
      \fi
      \draw[NNLine,line width=.5pt] ({11+6*\j},{-14-6*\i}) rectangle ++(6,6);
    }
  }
  \draw[NNHidden,line width=1.2pt] (17,-44) rectangle (53,-38);
  \foreach \i/\word in {1/猫,2/在,3/院,4/里,5/追,6/狗}{
    \node[token] (s\i) at ({62+17*\i},-18) {\word};
  }
  \node[nn transform,minimum width=30mm] (local) at (130,-43) {更新“追”};
  \foreach \j in {4,5,6}{\draw[nn operator path] (s\j.south)--(local.north);}
  \foreach \j in {1,2,3}{\node[nn annotation] at ({62+17*\j},-30) {$\times$};}
  \node[nn annotation] at (115,-55) {$36$个位置对 $\longrightarrow 16$个};
  \draw[FigureGrid,line width=.85pt] (0,-62)--(168,-62);

  \node[nn heading,anchor=west] at (0,-69) {(b) 线性：先汇总键值，再让查询读取};
  \node[anchor=north west,inner sep=0pt] at (0,-78) {\begin{tabular}{@{}ccc@{}}
    词元 & $\phi(\bm k_j)^{\top}$ & $\bm v_j^{\top}$\\
    猫 & $(1,0)$ & $(1,0)$\\
    在 & $(0,1)$ & $(0,1)$\\
    院 & $(1,1)$ & $(1,1)$\\
    里 & $(1,0)$ & $(2,0)$\\
    追 & $(0,1)$ & $(0,2)$\\
    狗 & $(1,1)$ & $(2,2)$
  \end{tabular}};
  \node[module,minimum width=37mm,align=center] (state) at (88,-94)
    {$\bm T=\begin{bmatrix}6&3\\3&6\end{bmatrix}$\\[2pt]$\bm b=(4,4)^{\top}$};
  \draw[nn operator path] (58,-94)--node[above] {求和} (state.west);
  \node[token,minimum width=40mm,align=center] (query) at (144,-82)
    {“追”的查询\\$\phi(\bm q)=(1,1)^{\top}$};
  \node[nn transform,minimum width=43mm,align=center] (linearout) at (144,-106)
    {$\bm z^{\top}=\dfrac{(1,1)\bm T}{(1,1)\bm b}$\\[2pt]$=(9/8,9/8)$};
  \draw[nn flow] (query.south)--(linearout.north);
  \draw[nn operator path] (state.east)--(115,-94)--(115,-106)--(linearout.west);
  \draw[FigureGrid,line width=.85pt] (0,-130)--(168,-130);

  \node[nn heading,anchor=west] at (0,-137) {(c) 低秩：把6个键值位置压成2个槽};
  \node[anchor=north west,inner sep=0pt] at (0,-146) {\begin{tabular}{@{}cccccc@{}}
    猫 & 在 & 院 & 里 & 追 & 狗\\
    $1/3$ & $1/3$ & $1/3$ & $0$ & $0$ & $0$\\[2pt]
    $0$ & $0$ & $0$ & $1/3$ & $1/3$ & $1/3$
  \end{tabular}};
  \node[nn annotation,anchor=west] at (0,-171) {$\bm E$};
  \node[module,minimum width=43mm,align=center] (slots) at (92,-157)
    {$\overline{\bm K}=\bm E\bm K$\\$\overline{\bm V}=\bm E\bm V$\\每行是一组混合表示};
  \draw[nn operator path] (58,-157)--(slots.west);
  \foreach \i/\word in {1/猫,2/在,3/院,4/里,5/追,6/狗}{
    \node at (137,{-142-5*\i}) {\word};
    \foreach \j in {1,2}{
      \draw[NNLine,fill=NNHidden!25,line width=.5pt]
        ({139+8*\j},{-144.5-5*\i}) rectangle ++(8,5);
    }
  }
  \node at (151,-141) {槽1};
  \node at (159,-141) {槽2};
  \draw[nn operator path] (slots.east)--(131,-157);
\end{tikzpicture}
